\chapter{Introduction}
\label{chap:chapter1}

This thesis is focused on the design, performance study and characterization of gaseous and hybrid X-ray detectors based on a custom, pixelated CMOS chip for applications in high-resolution astronomical polarimetry, spectroscopy and imaging. Specifically, this work combines experimental measurements on prototypes with the development of Monte Carlo simulations and dedicated event reconstruction algorithms.

The development of these detectors relies on the experience acquired in the last two decades during the realization of Gas Pixel Detectors (GPDs), a class of X-ray photoelectric polarimeters employed on board the Imaging X-ray Polarimetry Explorer (IXPE) mission funded by NASA and Agenzia Spaziale Italiana (ASI). IXPE, launched in 2021 and still operating, is dedicated to the study of polarization, combined with imaging, spectroscopy and timing of astrophysical sources in the energy range between 2 and 8~keV.

GPDs represented a breakthrough in X-ray polarimetry thanks to the use of a Gas Electron Multiplier (GEM) combined with a finely pixelated ASIC to track a photoelectron released after an X-ray photon is absorbed in the active gas volume. Tracking the photoelectron is the core working principle of a photoelectric polarimeter because, under specific conditions, the distribution of photoelectron emission angles depends on the source's degree of polarization. By accumulating enough statistics, it is possible to infer the polarization level of the source.

A fundamental component of the GPD is its custom ASIC, named XPOL, which performs an event-driven and self-triggering readout by dynamically defining a small region of interest around the event, significantly reducing the readout time compared to a full-frame acquisition. Despite this, current GPDs do not meet the speed requirements of the next generation of X-ray polarimeters, such as the one on board the enhanced X-ray Timing and Polarimetry (eXTP) mission, which will employ larger mirrors and will therefore need to be capable of detecting events with a much higher rate than IXPE. To address this, the XPOL electronics has been redesigned to increase the event readout speed and to define more efficiently and dynamically the region of interest. This new version has been developed in recent years and demonstrated to improve the readout speed by about an order of magnitude.

Motivated by these perspectives, this thesis builds upon the XPOL readout architecture to address two parallel goals: developing the next generation of GPDs by solving the physical limitations of the current multiplication stage, and extending this technology to solid-state sensors for non-polarimetric applications.

Regarding the first goal, current GPDs show some issues related to the GEM. Indeed, the GEM manufacturing process is characterized by imperfections that are the cause of variations in the detector azimuthal response (spurious modulation), which leads to the detection of a false polarization even when observing unpolarized sources. Another aspect of the GEM that has to be improved is the rate-dependent variation of the charge multiplication factor over time, caused by charge accumulation on the exposed dielectric inside the holes, leading to a degradation of the energy resolution. A possible solution to reduce these systematic effects and improve detector stability over time consists of replacing the GEM with micro-gap-like structures fabricated directly on top of the ASIC, creating a monolithic detector with the multiplication stage integrated on the readout chip. This concept defines the $\mu$GPD design. By fabricating the multiplication stage directly on top of the ASIC during the foundry manufacturing process, the tolerances are greatly reduced with respect to the GEM manufacturing process, reaching an accuracy on sub-micrometer length scale. Furthermore, with this design the surface of exposed dielectric of the multiplication structures can be much smaller compared to a GEM hole, significantly reducing the amount of charge that can be trapped during the avalanche and improving the stability of the detector gain over time. To achieve this goal, it is first necessary to demonstrate the feasibility of the concept by testing and characterizing the multiplication structures under realistic operating conditions before commissioning the final detector. To this end, prototype structures were fabricated via post-processing directly onto XPOL ASICs. In this context, my work is focused on the characterization of these test structures, evaluating their breakdown limits, gain capabilities, time stability and energy resolution.

In parallel with the gas detector developments, the fast readout capabilities and low-noise electronics of the XPOL ASICs have paved the way to the development of a new class of X-ray detectors for non-polarimetric applications requiring energy-resolved imaging, such as X-ray diffraction, material science and X-ray astrophysics. This is the Analog Spectral Imager for X-rays (ASIX) project, which consists of the development of Hybrid Pixel Detectors in which a solid-state sensor is coupled to an XPOL-family ASIC, inheriting GPDs' event-driven and self-triggering readout. The event-driven approach enables taking advantage of charge sharing, typical of detectors with small pixel pitches like the 50~$\mu$m hexagonal pixels of XPOL, to achieve better spatial resolution through offline analysis. At the same time, optimizing sensor design parameters that influence the charge sharing geometry allows for high energy resolution by maximizing the signal-to-noise ratio, which depends on the number of pixels involved in the event. My work in the ASIX project was focused on the development of the Monte Carlo simulation and analysis software, the calibration of the detector and the development of different algorithms for offline reconstruction of incidence position and energy of the photons. These algorithms were evaluated using both dedicated simulations and experimental data collected with a prototype. Through these analyses, the spatial and energy resolution of the detector were determined. Finally, a comparison between simulations and laboratory measurements was carried out to verify the validity of the Monte Carlo simulation software.

This thesis is structured as follows:

Following this introductory Chapter 1, Chapter 2 provides a short review of X-ray detection in astronomy, covering the historical evolution of space observatories alongside the fundamental light-matter interaction mechanisms and the operating principles of gas and solid-state detectors.  

Chapter 3 focuses on the IXPE mission and the GPD, describing its architecture, the technological limitations that introduce systematic effects and limit the readout speed, and the XPOL-III upgrade of the readout ASIC.  

Chapter 4 presents the development of $\mu$GPDs, evaluating the experimental performance of these prototypes across different stages of development. The performance is evaluated in terms of gain, energy resolution, and temporal stability.  

Chapter 5 is dedicated to the ASIX project, describing the detector design and the Monte Carlo simulation and analysis framework developed for calibration and event reconstruction. Offline algorithms for photon position and energy reconstruction are presented and validated using simulations and laboratory data.